\documentclass[10pt]{article}

% Specify the margins and text size.
\setlength{\textwidth}{6.5in}
\setlength{\textheight}{9.5in}
\setlength{\oddsidemargin}{0pt}
\setlength{\evensidemargin}{0pt}
\setlength{\topmargin}{0pt}
\setlength{\hoffset}{-.05in}
\setlength{\voffset}{-.8in}

\setlength{\parskip}{5pt}
\setlength{\parindent}{0pt}

% Load some fonts and symbol packages
\usepackage{latexsym}
\usepackage{pifont}       % contains 'star' symbol for counterinsurgency handbook title
\usepackage{yfonts} 
\usepackage{amsmath}
\usepackage{amsfonts}

\usepackage{graphicx}     % actually, this is loaded by pstricks
\usepackage[T1]{fontenc}
\usepackage{ifthen}
\usepackage{pstricks,pst-grad,pst-text,pst-node,multido,pst-plot,calc,pst-3dplot}
%\usepackage[all]{xy}
%\usepackage{animate}

% The hyperref package inserts links into the pdf.
\definecolor{MyLinkColor}{rgb}{.1,.2,1}
\definecolor{MyCiteColor}{rgb}{.1,1,.2}
\definecolor{MyURLColor}{rgb}{.4,.4,.4}
\usepackage[backref=true,pagebackref=false,hyperindex,colorlinks=true,
  linkcolor=MyLinkColor,urlcolor=MyURLColor]{hyperref}


% The tweaklist package is something I found on the web.  It provides a simple interface
% for making changes to spacing used in the itemize and enumerate environments.  Comment
% this out if you don't care to use tweaklist.
\usepackage{tweaklist}
\renewcommand{\itemhook}{\setlength{\parskip}{2pt}\setlength{\parsep}%
{1pt}\setlength{\topsep}{0pt}\setlength{\itemsep}{0pt}}

\newcommand{\U}{\underline{\hspace{5pt}}}

\usepackage{listings}
\newcounter{EX}\setcounter{EX}{1}
\newcommand{\EXERCISE}{\arabic{EX}.\stepcounter{EX} }

\begin{document}
\pagestyle{empty}
\lstset{language=R, showspaces=false, showstringspaces=false}

\href{http://www.su.edu}{\includegraphics[height=1.75cm]{sulogo.eps}}
\vspace{-1.69cm}

{\small \hfill
\begin{tabular}{cl}
& Math 207\\
& Introduction to  Statistics\\
\end{tabular}
}
\setlength{\baselineskip}{1.05\baselineskip}

\begin{center}
\textbf{\large  Regression Review Exercises}
\end{center}

1. Here is some data:
\begin{center}
\begin{tabular}{cc}
$x$ & $y$\\\hline
0   & 1\\
1   & 3\\
2   & 2
\end{tabular}
\end{center}

\hspace{10pt} a) Sketch the data and the line $y=x+1$ then calculate the RMS error for the line 
\vspace{1.4in}

\hspace{10pt} b) Sketch the data and the line $y=\frac{3}{2}\,x+1$ 
then calculate the RMS error for the line 
\vspace{1.3in}

\hspace{10pt} c) Sketch the data and the line $y=\frac{2}{3}\,x+1$ 
then calculate the RMS error for the line 
\vspace{1.3in}

\hspace{10pt} d) For the given data, $\mbox{sd}_x=\mbox{sd}_y=\sqrt{2/3}$ and $r=1/2$.
Calculate
the {sd} line and find its RMS error.
\vspace{1.4in}

\hspace{10pt} e) Calculate the regression line and find its RMS error.  
\vfill
\eject

2.  For a given data set we find that
$\mbox{mean}(x)=10$, $\mbox{mean}(y)=0$, ${sd}_x=1$, ${sd}_y=2$ and  $r=1/5$.

\hspace{10pt} a) Predict $y$ if $x=11$.
\vspace{1.2in}

\hspace{10pt} b)  If $x=11$, about what range of $y$ values contains about 68\% of the data?
\vspace{.5in}

\hspace{10pt} c)  If $x=11$, about what range of $y$ values contains about 95\% of the data?
\vspace{.5in}

\hspace{10pt} d)  If $x=8$, about what range of $y$ values contains about 68\% of the data?
\vspace{.5in}

\hspace{10pt} e)  If $x=8$, about what range of $y$ values contains about 95\% of the data?
\vspace{.5in}


\hspace{10pt} f) Predict $x$ if $y=2$.
\vspace{1.2in}

\hspace{10pt} g)  If $y=2$, about what range of $x$ values contains about 68\% of the data?
\vspace{.5in}

\hspace{10pt} h)  If $y=2$, about what range of $x$ values contains about 95\% of the data?
\vspace{.5in}

\hspace{10pt} i)  If $y=-1$, about what range of $x$ values contains about 68\% of the data?
\vspace{.5in}

\hspace{10pt} j)  If $y=-1$, about what range of $x$ values contains about 95\% of the data?
\vfill
\eject

3. Calculate the given area under the standard normal curve.

\hspace{10pt} a) to the left of $z=1.5$
\vspace{.5in}

\hspace{10pt} b) to the right of $z=2.5$
\vspace{.5in}

\hspace{10pt} c) between $-1.5$ and $2.5$
\vspace{1in}

4. For what value of $z$ is each statement true?

\hspace{10pt} a) area to the left of $z$ is 35\%
\vspace{.5in}

\hspace{10pt} b) area to the right of $z$ is 10\%
\vspace{.5in}

\hspace{10pt} c) area between $z$ and $-z$  is 10\%
\vspace{.5in}

5. Suppose a measurement has a normal distribution with $\mbox{mean}=10$ and $\mbox{{sd}}=4$.
Calculate the percent of measurements that satisfy the following.

\hspace{10pt} a) at most 16
\vspace{1in}

\hspace{10pt} b) at least 0
\vspace{1in}

\hspace{10pt} c) between 4 and 20
\vfill
\eject


\end{document}

